% GENERATED FILE. Edit canonical lessons, then run npm run generate:books.
% canonical-source-hash: fnv1a64:1bf2c3fd3ea5ae27
% canonical-lessons: PT-C199-milho, PT-C199-canela, PT-C199-ementa, PT-C199-aquecimento, PT-C199-torneira

\chapter{Maize, Cinnamon, Menu, Heating, Tap}
\label{ch:pt-maize-cinnamon-menu-heating-tap}
\hlchaptermodality{199}

\begin{chapteropening}
\textbf{By the end of this chapter:} \emph{I can say maize, cinnamon, a menu, heating and a tap.}

Complete the last lesson of chapter 199: I can say maize, cinnamon, a menu, heating and a tap.
\end{chapteropening}

\section[o milho]{o milho — maize, sweetcorn}

\label{lesson:PT-C199-milho}



\begin{quote}
Before the new one: say the Portuguese for a packet, then the Portuguese for a tin.
\end{quote}

\subsection*{You'll want to know: o milho}
\textbf{o milho} — \textquotedblleft{}maize, sweetcorn\textquotedblright{}.

\begin{cousinweb}
\textbf{milho} $\leftarrow$ Latin \emph{milium}, \textquotedblleft{}millet\textquotedblright{} (EtymDB 2.1) — inherited. Maize from America took the name of an old grain.
\end{cousinweb}

\begin{grammarlens}[title={What you've built}]
One more word for this chapter.
\end{grammarlens}

\subsection*{Your turn}
\begin{itemize}
\raggedright
  \item \emph{Say it:} \emph{o milho}
  \item \emph{Say it:} \emph{o milho}, once more
  \item \emph{Say it:} say \emph{o milho}
  \item \emph{Recall:} say the Portuguese for a packet, then the Portuguese for a tin, then say \emph{o milho} again
\end{itemize}

\begin{tcolorbox}[breakable,colback=teal!4,colframe=teal!35!black,title={Before you move on}]
What does o milho mean? (\textquotedblleft{}Maize, sweetcorn\textquotedblright{}.) Say it once more.
\end{tcolorbox}
\section[a canela]{a canela — cinnamon}

\label{lesson:PT-C199-canela}



\begin{quote}
Before the new one: say the Portuguese for a tin, then the Portuguese for maize.
\end{quote}

\subsection*{You'll want to know: a canela}
\textbf{a canela} — \textquotedblleft{}cinnamon\textquotedblright{}.

\begin{cousinweb}
\textbf{canela} $\leftarrow$ Medieval Latin \emph{canella}, \textquotedblleft{}little reed\textquotedblright{} (EtymDB 2.1), from \emph{canna}, \textquotedblleft{}reed\textquotedblright{} — a learned borrowing from Medieval Latin. Rolled cinnamon bark looks like little reeds.
\end{cousinweb}

\begin{grammarlens}[title={What you've built}]
One more word for this chapter.
\end{grammarlens}

\subsection*{Your turn}
\begin{itemize}
\raggedright
  \item \emph{Say it:} \emph{a canela}
  \item \emph{Say it:} \emph{a canela}, once more
  \item \emph{Say it:} say \emph{a canela}
  \item \emph{Recall:} say the Portuguese for a tin, then the Portuguese for maize, then say \emph{a canela} again
\end{itemize}

\begin{tcolorbox}[breakable,colback=teal!4,colframe=teal!35!black,title={Before you move on}]
What does a canela mean? (\textquotedblleft{}Cinnamon\textquotedblright{}.) Say it once more.
\end{tcolorbox}
\section[a ementa]{a ementa — a menu}

\label{lesson:PT-C199-ementa}



\begin{quote}
Before the new one: say the Portuguese for maize, then the Portuguese for cinnamon.
\end{quote}

\subsection*{You'll want to know: a ementa}
\textbf{a ementa} — \textquotedblleft{}a menu\textquotedblright{}.

\begin{cousinweb}
\textbf{ementa} — its origin is uncertain. EtymDB 2.1 links it only to a verb meaning \textquotedblleft{}to list\textquotedblright{}.
\end{cousinweb}

\begin{grammarlens}[title={What you've built}]
One more word for this chapter.
\end{grammarlens}

\subsection*{Your turn}
\begin{itemize}
\raggedright
  \item \emph{Say it:} \emph{a ementa}
  \item \emph{Say it:} \emph{a ementa}, once more
  \item \emph{Say it:} say \emph{a ementa}
  \item \emph{Recall:} say the Portuguese for maize, then the Portuguese for cinnamon, then say \emph{a ementa} again
\end{itemize}

\begin{tcolorbox}[breakable,colback=teal!4,colframe=teal!35!black,title={Before you move on}]
What does a ementa mean? (\textquotedblleft{}A menu\textquotedblright{}.) Say it once more.
\end{tcolorbox}
\section[o aquecimento]{o aquecimento — heating}

\label{lesson:PT-C199-aquecimento}



\begin{quote}
Before the new one: say the Portuguese for cinnamon, then the Portuguese for a menu.
\end{quote}

\subsection*{You'll want to know: o aquecimento}
\textbf{o aquecimento} — \textquotedblleft{}heating\textquotedblright{}.

\begin{cousinweb}
\textbf{aquecimento} is built from the verb \textquotedblleft{}to heat\textquotedblright{} and \emph{-mento}. That verb goes back to Latin \emph{calescere}, \textquotedblleft{}to grow warm\textquotedblright{} (Lewis \& Short).
\end{cousinweb}

\begin{grammarlens}[title={What you've built}]
One more word for this chapter.
\end{grammarlens}

\subsection*{Your turn}
\begin{itemize}
\raggedright
  \item \emph{Say it:} \emph{o aquecimento}
  \item \emph{Say it:} \emph{o aquecimento}, once more
  \item \emph{Say it:} say \emph{o aquecimento}
  \item \emph{Recall:} say the Portuguese for cinnamon, then the Portuguese for a menu, then say \emph{o aquecimento} again
\end{itemize}

\begin{tcolorbox}[breakable,colback=teal!4,colframe=teal!35!black,title={Before you move on}]
What does o aquecimento mean? (\textquotedblleft{}Heating\textquotedblright{}.) Say it once more.
\end{tcolorbox}
\section[a torneira]{a torneira — a tap}

\label{lesson:PT-C199-torneira}



\begin{quote}
Before the new one: say the Portuguese for a menu, then the Portuguese for heating.
\end{quote}

\subsection*{You'll want to know: a torneira}
\textbf{a torneira} — \textquotedblleft{}a tap\textquotedblright{}.

\begin{cousinweb}
\textbf{torneira} is built on the word for a turn ($\leftarrow$ Latin \emph{tornare}, \textquotedblleft{}to turn on a lathe\textquotedblright{}) and \emph{-eira}: the thing you turn.
\end{cousinweb}

\begin{grammarlens}[title={What you've built}]
One more word for this chapter.
\end{grammarlens}

\subsection*{Your turn}
\begin{itemize}
\raggedright
  \item \emph{Say it:} \emph{a torneira}
  \item \emph{Say it:} \emph{a torneira}, once more
  \item \emph{Say it:} say \emph{a torneira}
  \item \emph{Recall:} say the Portuguese for a menu, then the Portuguese for heating, then say \emph{a torneira} again
\end{itemize}

\begin{tcolorbox}[breakable,colback=teal!4,colframe=teal!35!black,title={Before you move on}]
What does a torneira mean? (\textquotedblleft{}A tap\textquotedblright{}.) Say it once more.
\end{tcolorbox}
